\documentclass[10pt]{article}
\usepackage[a4paper,landscape,margin=1.5cm]{geometry}
\usepackage{fontspec}
\setmainfont{DejaVu Serif}
\setmonofont{DejaVu Sans Mono}[Scale=0.9]
\usepackage{amsmath,amssymb,longtable,booktabs,array,xurl}
\usepackage[hidelinks]{hyperref}
\renewcommand{\arraystretch}{1.25}
\setlength{\tabcolsep}{4pt}
\setlength{\parindent}{0pt}
\setlength{\parskip}{4pt}
\title{Inequality and Social Contests: measurement map}
\date{Pass 2}
\begin{document}
\maketitle
Micro model only. Macro variables deferred. Model claims M1--M23 are those of the manuscript skeleton of the same pass. Each row says how an observer could obtain the object. Each observation also carries its support. PROVEN means a Lean theorem shows the observed quantity determines the object within the model; DEFINITIONAL means the observation is the object once the referent is accepted; CONSTRUCT means the object is not measured; UNVERIFIED means neither, and no headline may rest on it. Entries marked PROPOSED await the author's confirmation; a headline resting on one is flagged, and a headline resting on a row pending on an undecided model claim is an error.

\section*{Support vocabulary}
\begin{description}
\item[\texttt{LITERATURE}] a literature claim uses the variable with this real-world meaning
\item[\texttt{POSIT}] the paper's own identification, supported by nothing else in the ledger
\item[\texttt{NONE}] modelling device with no real-world referent

\section*{Observation support}

\item[\texttt{PROVEN}] a Lean theorem shows the observed quantity determines the object within the model
\item[\texttt{DEFINITIONAL}] the observed quantity is the object itself once the referent is accepted
\item[\texttt{LITERATURE}] a literature row read in full measures the object
\item[\texttt{CONSTRUCT}] not measured; a representation of recorded choices
\item[\texttt{UNVERIFIED}] neither proven nor measured in a paper read in full; no headline may rest on it

\section*{Observation status}

\item[\texttt{PROPOSED}] how an observer could obtain the object, proposed and awaiting the author's confirmation
\item[\texttt{CONFIRMED}] confirmed by the author
\end{description}

\section{Inputs}
{\small
\begin{longtable}{>{\raggedright\arraybackslash}p{0.8cm} >{\raggedright\arraybackslash}p{2.2cm} >{\raggedright\arraybackslash}p{4.0cm} >{\raggedright\arraybackslash}p{3.6cm} >{\raggedright\arraybackslash}p{1.6cm} >{\raggedright\arraybackslash}p{4.0cm} >{\raggedright\arraybackslash}p{4.6cm} >{\raggedright\arraybackslash}p{1.8cm}}
\toprule ID & Symbol & Model meaning & Real-world referent & Support & Direction & How an observer could obtain it & Used in \\ \midrule \endhead
X1 & $N$ & number of workers, equal to the number of jobs & size of the labour market & \texttt{POSIT} & --- & count of jobs in the market studied, from a job-title register or a matched employer-employee panel \newline \texttt{PROPOSED} \newline \texttt{DEFINITIONAL}. The observation is the object itself once the referent is accepted & M1, M10, M23 \\ \midrule
X2 & $q\in[0,1)$ & rank quantile of a job on the ladder & percentile of a job in the wage ranking & \texttt{POSIT} & higher $q$ is a better-paid job (M2) & percentile of the job's wage in the wage distribution of the same market, from payroll or matched employer-employee data \newline \texttt{PROPOSED} \newline \texttt{PROVEN} (\nolinkurl{Isc.quantile_from_wage}, \nolinkurl{Isc.pareto_strictMono_quantile}). On the ladder the wage determines the quantile, $q=1-(x/x_0)^{-a}$, and the wage is strictly increasing in $q$ & M2, M3, M4 \\ \midrule
X3 & $x_0$ & income at $q=0$ & lowest wage & \texttt{POSIT} & --- & lowest wage of the market, read as a low percentile of the wage distribution in payroll data \newline \texttt{PROPOSED} \newline \texttt{PROVEN} (\nolinkurl{Isc.pareto_at_zero}). $x(0)=x_0$ & M1, M2, M3, M4 \\ \midrule
X4 & $a$ & tail index of the income ladder $x(q)=x_0(1-q)^{-1/a}$ & income dispersion across ranks, inversely & \texttt{POSIT} & higher $a$ lowers income at every interior quantile (M3) and narrows every income gap (M4), so higher $a$ means a more equal ladder & Pareto index of the wage distribution, estimated as minus the slope of the log share of workers above a wage on the log of that wage \newline \texttt{PROPOSED} \newline \texttt{PROVEN} (\nolinkurl{Isc.log_survival_linear}). On the ladder $\log(1-q)=-a(\log x-\log x_0)$, so the slope of log survival share on log wage is $-a$ & M1, M2, M3, M4, M5, M21 \\ \midrule
X5 & $b$ & exponent of the upward-move cost $C(m_i,m_k)=(m_k-m_i)^b$ & convexity of retraining cost in the income distance climbed & \texttt{POSIT} & for $b\ge1$ two adjacent steps cost no more than the jump (M8). For $0\le b\le1$ the jump costs no more than the steps (M9) & growth of retraining cost (tuition and earnings forgone while retraining) with the wage gap between origin and destination job, from training records linked to job moves \newline \texttt{PROPOSED} \newline \texttt{PROVEN} (\nolinkurl{Isc.cost_exponent_log_linear}). $\log C=b\log(m_k-m_i)$, so the slope of log cost on log wage gap is $b$ & M6, M7, M8, M9, M21, M23 \\ \midrule
X6 & $m$ & income of a position, written $x$ on the ladder & wage of a job & \texttt{POSIT} & --- & wage of the job, from payroll data \newline \texttt{PROPOSED} \newline \texttt{DEFINITIONAL}. The observation is the object itself once the referent is accepted & M6, M7, M8, M9, M18 \\ \midrule
X7 & $L$ & fixed cost of entering the lottery & cost of social navigation, such as networking, club membership and compliance with social closure & \texttt{POSIT} & higher $L$ lowers the participation threshold (M13) & spending on social navigation such as membership fees and networking time, from expenditure or time-use surveys \newline \texttt{PROPOSED} \newline \texttt{UNVERIFIED}. That spending on networking measures the entry cost is a posit; no paper read in full measures it & M11, M12, M13, M20, M24, M25, M26, M27 \\ \midrule
X8 & swap rule & a lottery entrant $i$ swaps positions with a partner $k$ drawn uniformly from the other entrants (M14, decided 10 Oct 2026); a draw across status groups is weighted by $1-\sigma$ (M29) & a social contest that reallocates jobs by connection or luck rather than by skill & \texttt{POSIT} & --- & not observed directly; its implication is a rank change after job loss that is unrelated to the worker's own rank, testable on the distribution of rank changes in panel data \newline \texttt{PROPOSED} \newline \texttt{UNVERIFIED}. The implication for rank changes is not proven, and no paper read in full measures it & M10, M11, M14, M24, M29 \\ \midrule
X9 & $u$ & utility over income in expected-utility form, unique up to a positive affine transformation (decision of 10 Oct 2026), S-shaped in the manuscript, functional form not given & attitude to income risk, risk-seeking below a reference income and risk-averse above it & \texttt{POSIT} & --- & not measured; a construct that rationalises recorded choices between the lottery route and the downgrade at each income \newline \texttt{PROPOSED} \newline \texttt{CONSTRUCT}. Not measured, by the decision of 10 Oct 2026 (Q4); a representation of recorded choices & M15, M16, M18, M19, M20 \\ \midrule
X10 & $\bar y$ & mean of net income from the lottery component & expected income & \texttt{POSIT} & --- & mean income change of workers who take the lottery route, from panel data \newline \texttt{PROPOSED} \newline \texttt{DEFINITIONAL}. The sample mean estimates the mean income change; the object is that mean once the referent is accepted & M15 \\ \midrule
X11 & $\sigma_y$ & standard deviation of net income from the lottery component & income risk created by taking part in social contests & \texttt{POSIT} & --- & standard deviation of that income change, from panel data \newline \texttt{PROPOSED} \newline \texttt{DEFINITIONAL}. The sample standard deviation estimates it once the referent is accepted & M15, M16, M17, M18 \\ \midrule
X12 & $\varepsilon$ & standardised shock, $\mathbb{E}\varepsilon=0$, $\mathbb{E}\varepsilon^2=1$ & --- & \texttt{NONE} & --- & --- & M15 \\ \midrule
X13 & $s$ & scale in $\sigma_y=s(1-\lambda)m$ & income risk per unit of income staked on social contests & \texttt{POSIT} & --- & ratio of the income risk of the lottery route to the income staked, compared across income levels in panel data \newline \texttt{PROPOSED} \newline \texttt{PROVEN} (\nolinkurl{Isc.risk_scale_identified}). $s=\sigma_y/((1-\lambda)m)$ whenever $(1-\lambda)m\ne0$ & M18 \\ \midrule
X14 & $\lambda$ & share of income allocated to upskilling, with $1-\lambda$ allocated to the lottery & share of resources spent on human capital rather than on social navigation & \texttt{POSIT} & --- & share of spending on training or certification against spending on social navigation, from expenditure or time-use surveys \newline \texttt{PROPOSED} \newline \texttt{UNVERIFIED}. That the spending shares measure the upskilling share is a posit; no paper read in full measures it & M17, M18, M22 \\ \midrule
X22 & $S_P$ & sum of the incomes of the lottery entrants & income held by those who compete through connections & \texttt{POSIT} & --- & sum of the wages of workers who take the lottery route, from panel data \newline \texttt{PROPOSED} \newline \texttt{DEFINITIONAL}. The observation is the object itself once the referent is accepted & M11, M13, M24 \\ \midrule
X23 & $|P|$ & number of lottery entrants & number of workers competing through connections & \texttt{POSIT} & a risk-neutral pool of two or more with $L>0$ is not stable (M24) & count of workers who take the lottery route, from panel data \newline \texttt{PROPOSED} \newline \texttt{DEFINITIONAL}. The observation is the object itself once the referent is accepted & M11, M12, M13, M24 \\ \midrule
X24 & $h$ & categorical status label of a worker, one of two groups, fixed within a working life (M28) & status group in Weber's sense, a categorical identity such as ethnicity, religion, nationality or political identity & \texttt{POSIT} & --- & group membership recorded in survey or administrative data, as self-reported identity \newline \texttt{PROPOSED} \newline \texttt{UNVERIFIED}. That a recorded identity is the status group of the model is a posit; no paper read in full measures it & M25, M26, M27, M28, M29 \\ \midrule
X25 & $\sigma$ & closure parameter in $[0,1]$; a partner draw across groups has weight $1-\sigma$ against $1$ within the group (M29) & social closure, the degree to which connections across status groups are restricted & \texttt{POSIT} & at $\sigma=0$ the contest ignores groups (M25), at $\sigma=1$ it runs within groups (M26); raising it lowers the gain of an entrant whose out-group is relatively richer (M27) & share of lottery-route moves whose partner is from the other group, with the numbers of entrants in each group, from referral or network data linked to a matched employer-employee panel \newline \texttt{PROPOSED} \newline \texttt{PROVEN} (\nolinkurl{Isc.crossShare_strictAnti}, \nolinkurl{Isc.closure_from_crossShare}). Under the draw rule of M29, with $n$ in-group and $m$ out-group entrants besides the worker, the share of cross-group partners is $(1-\sigma)m/(n+(1-\sigma)m)$, strictly decreasing in $\sigma$, and $\sigma=1-nc/(m(1-c))$ for an observed share $c$; it needs out-group entrants, $m>0$ & M25, M26, M27, M29 \\ \midrule
\end{longtable}}

\section{Outputs}
{\small
\begin{longtable}{>{\raggedright\arraybackslash}p{0.8cm} >{\raggedright\arraybackslash}p{2.2cm} >{\raggedright\arraybackslash}p{4.0cm} >{\raggedright\arraybackslash}p{3.6cm} >{\raggedright\arraybackslash}p{1.6cm} >{\raggedright\arraybackslash}p{4.0cm} >{\raggedright\arraybackslash}p{4.6cm} >{\raggedright\arraybackslash}p{1.8cm}}
\toprule ID & Symbol & Model meaning & Real-world referent & Support & Direction & How an observer could obtain it & Used in \\ \midrule \endhead
X15 & $x(q)$\newline produced by M2 & $x_0(1-q)^{-1/a}$ & wage of the job at percentile $q$ & \texttt{POSIT} & increasing in $q$ (M2), decreasing in $a$ (M3) & wage at percentile q of the wage distribution, from payroll data \newline \texttt{PROPOSED} \newline \texttt{PROVEN} (\nolinkurl{Isc.quantile_from_wage}). The wage at percentile $q$ is $x(q)$, and $q$ is recovered from it & M2, M3, M4, M10, M11, M13 \\ \midrule
X16 & $C(m_i,m_k)$\newline produced by M7 & $(m_k-m_i)^b$ if $m_i<m_k$, $0$ otherwise & retraining cost of moving from job $i$ up to job $k$ & \texttt{POSIT} & zero for downward and lateral moves (M7) & retraining cost of an observed upward move, from training records and earnings forgone \newline \texttt{PROPOSED} \newline \texttt{DEFINITIONAL}. The observation is the object itself once the referent is accepted & M6, M7, M8, M9 \\ \midrule
X17 & $S$\newline produced by M10 & $\sum_k x_k$ & total wage bill & \texttt{POSIT} & unchanged by any swap (M10) & total wage bill of the market, from payroll data \newline \texttt{PROPOSED} \newline \texttt{DEFINITIONAL}. The observation is the object itself once the referent is accepted & M10 \\ \midrule
X18 & $G_i$\newline produced by M11 & $\frac{S_P-|P|x_i}{|P|-1}-L$ & expected net payoff to worker $i$ from social navigation, valued risk-neutrally & \texttt{POSIT} & sums to $-|P|L$ over the pool (M12); negative for the richest entrant when $L>0$ (M24) & mean income change net of the entry cost for workers from rank i who take the lottery route, from panel data \newline \texttt{PROPOSED} \newline \texttt{PROVEN} (\nolinkurl{Isc.poolGain_closed}). The expected income change of an entrant under a uniform draw is the average of $x_k-x_i$ over the other entrants, which M11 equates to $G_i+L$ & M11, M12, M13, M24, M25, M26 \\ \midrule
X19 & $\bar x_{\mathrm{part}}$\newline produced by M13 & $\frac{S_P-(|P|-1)L}{|P|}$, with $G_i>0\iff x_i<\bar x_{\mathrm{part}}$ & income below which social navigation pays in expectation & \texttt{POSIT} & --- & income below which workers take the lottery route, estimated as a participation threshold on rank \newline \texttt{PROPOSED} \newline \texttt{UNVERIFIED}. M13 identifies the threshold only under risk-neutral valuation, and M24 shows a risk-neutral pool of two or more is not stable, so an observed threshold also reflects the curvature of $u$ & M13 \\ \midrule
X20 & $\tfrac12u^{\prime\prime}(\bar y)\sigma_y^2$\newline produced by M15 & second-order risk term in expected utility & utility value of the income risk from social contests & \texttt{POSIT} & positive exactly where $u^{\prime\prime}>0$ (M16) & not measured; a construct inferred only through choices under risk \newline \texttt{PROPOSED} \newline \texttt{CONSTRUCT}. Not measured, by the decision of 10 Oct 2026 (Q4); a representation of recorded choices & M15, M16 \\ \midrule
X21 & $p$\newline produced by M18 & $\tfrac12u^{\prime\prime}s^2m^2$, coefficient on $(1-\lambda)^2$ & risk-premium coefficient of a worker with income $m$ & \texttt{POSIT} & scales with $m^2$ (M18) & not measured directly; inferred from how take-up of the lottery route falls with income \newline \texttt{PROPOSED} \newline \texttt{CONSTRUCT}. Inferred from choices, never measured directly (decision of 10 Oct 2026, Q4) & M18 \\ \midrule
X26 & $G_i(\sigma)$\newline produced by M25, M26, M27 & $\frac{a+(1-\sigma)b}{n+(1-\sigma)m}-L$, the closure-weighted expected swap gain & expected net payoff from competing through connections when connections are closed along status lines & \texttt{POSIT} & equals $G_i$ of M11 at $\sigma=0$ (M25) & mean income change net of the entry cost for entrants of each group who take the lottery route, from panel data \newline \texttt{PROPOSED} \newline \texttt{PROVEN} (\nolinkurl{Isc.closureGain_formula}). The closure-weighted gain is the weighted average of $x_k-x_i$ over the other entrants, weights $1$ within and $1-\sigma$ across groups & M25, M26, M27 \\ \midrule
\end{longtable}}
\end{document}
